\newpage

\subsection{Übertragungsverhalten von Verstärkerschaltungen}\label{sec:1-4}
\subsubsection{Aufgabenstellung}

In diesem Abschnitt wird das Übertragungsverhalten eines invertierenden und eines nicht
invertierenden Operationsverstärkers untersucht. Dazu wird der OPV \texttt{A3} (LM741)
mit dem jeweiligen Widerstandsnetzwerk gegenkoppelt, wobei in beiden Fällen eine
Verstärkung von ca. \SI{40}{\decibel} eingestellt wird. Es wird untersucht, wie die 
Ausgangsspannung $U_A$ auf sinusförmige Eingangsspannungen $U_E$ bei verscheidenen
Frequenzen reagiert. 

\subsubsection{Verwendetes Material}

\begin{table}[h]
	\centering
	\caption{Verwendetes Material}
	\begin{tabularx}{\textwidth}{XXX}
		\textbf{Material} & \textbf{Verwendung} & \textbf{Anmerkungen}\\ \midrule
		Messverstärker & $A_3, R_1, R_2$ & $A_{GK} \approx\SI{40}{\decibel}$\\
		Funktionsgenerator & Eingangsspannung $U_E$ & \\
		Oszilloskop & Ausgangsspannung $U_A$& \\
	\end{tabularx}	
\end{table}
\subsubsection{Versuchsaufbau}

\begin{figure}[h]
	\begin{minipage}[t]{.49\textwidth}
	\centering
	\begin{circuitikz}[scale=0.9, transform shape]
		\myIA{A3}{(0,0)}{$R_1$}{$R_2$};
		\draw (A3-in) to[open, v=$U_E$, o-o] (A3-in |- A3-gnd) node[ground]{};
		\draw (A3-out) to[open, v=$U_A$, o-o] (A3-out |- A3-gnd) node[ground]{};
	\end{circuitikz}
	\caption{Messschaltung i. Verstärker}\label{fig:i_amplifier}
	\end{minipage}\hfill
	\begin{minipage}[t]{.49\textwidth}
	\centering
	\begin{circuitikz}[scale=0.9, transform shape]
		\myNIA{A3}{(0,0)}{$R_1$}{$R_2$};
		\draw (A3-in) to[open, v=$U_E$, o-o] (A3-in |- A3-gnd) node[ground]{};
		\draw (A3-out) to[open, v=$U_A$, o-o] (A3-out |- A3-gnd) node[ground]{};
	\end{circuitikz}
	\caption{Messschaltung n.i. Verstärker}\label{fig:ni_amplifier}
	\end{minipage}
\end{figure}

Zu beachten ist, dass die Offsetspannung des OPV \texttt{A3} wie in
\autoref{sec:1-1} abgeglichen werden muss.

\subsubsection{Berechnungen}

Mittels Gegenkopplung wird eine Verstärkung von $A_{GK}\approx\SI{40}{\decibel}
= 10^{\frac{40}{20}} = 100$ eingestellt. Um diese zu erreichen, wurden die
Widerstandswerte wie folgt gewählt:

\begin{itemize}
	\item Invertierender Verstärker:
	\begin{equation*}
		A_{GK} = \frac{R_2}{R_1} = 100 \implies R_2 = 100 R_1 \implies \underline{R_1 = \SI{1}{\kilo\ohm}, R_2 = \SI{100}{\kilo\ohm}}
	\end{equation*}
	\item Nicht invertierender Verstärker:
	\begin{equation*}
		A_{GK} = 1 + \frac{R_2}{R_1} = 100 \implies R_2 = 99 R_1 \implies \underline{R_1 = \SI{1}{\kilo\ohm}, R_2 \approx \SI{100}{\kilo\ohm}}
	\end{equation*}
\end{itemize}

Für den i.-Verstärker wird das Verstärkungsziel genau erreicht, für den n.i.-Verstärker wurde ein Widerstand von $\SI{100}{\kilo\ohm}$ verwendet, da kein $\SI{99}{\kilo\ohm}$ Widerstand verfügbar war. Dies führt zu einer minimal höheren Verstärkung von $A_{GK} = 1 + \frac{\SI{100}{\kilo\ohm}}{\SI{1}{\kilo\ohm}} = 101$.

Auf Basis dieser Verstärkung wird die Spitzenspannung $U_{E,sp}$ so gewählt,
dass sich das Ausgangssignal bei der größten Verstärkung noch im linearen
Aussteuerbereich befindet. Dieser Bereich befindet sich jeweils ca.
\SI{3}{\volt} unterhalb der Versorgungsspannung \cite[S.~66]{EMTPraktikum2025}
von $\pm \SI{12}{\volt}$, also bei $\pm \SI{9}{\volt}$. Für beide
Verstärkergilt für die Wahl von
$U_{E,sp}=\SI{50}{\milli\volt}$:
\begin{equation*}
	U_{A,pp} = A_{GK} \cdot U_{E,pp} = 100 \cdot \SI{50}{\milli\volt} = \SI{5}{\volt} < \SI{9}{\volt}
\end{equation*}

\subsubsection{Durchführung}

\begin{wrapfigure}{r}{0.48\textwidth}
	\centering\vspace{-10pt}
	\includegraphics[width=0.48\textwidth]{BodeMessung.png}
	\caption{Messung der Dämpfung und Phasenverschiebung}
\end{wrapfigure}

Um nun die Übertragungsfunktion zu bestimmen, wird die Frequenz des
Eingangssignals von $\SI{10}{\hertz}$ bis $\SI{1}{\mega\hertz}$ in
logarithmischen gleichmäßigen abständen variiert. Die Messung der Phase erfolgt
durch die Messung der Zeitdifferenz $\Delta t$ zwischen Ein- und
Ausgangswellenberge, welche in eine Phasenverschiebung umgerechnet werden kann:
\begin{equation*}
	\varphi = \Delta t\cdot f \cdot \SI{360}{\degree}
\end{equation*}

Für die Messung des Betragsganges wird die Spitzenspannung des Ausgangssignals abgelesen und in Dezibel umgerechnet:
\begin{equation*}
	A_{dB} = 20 \cdot \log_{10}\left(\frac{U_{A,sp}}{U_{E,sp}}\right)
\end{equation*}

\newpage
\subsubsection{Messergebnisse}

Die Messergebnisse sind in \autoref{tab:i_bode} für den invertierenden
Verstärker und in \autoref{tab:ni_bode} für den nicht invertierenden Verstärker
zusammengefasst. Für den invertierenden Verstärker wurde die Frequenz nur bis
$\SI{800}{\kilo\hertz}$ statt $\SI{1}{\mega\hertz}$ erhöht, da das
Ausgangssignal darüber für eine aussagekräftige Messung zu stark verrauscht war.

\begin{table}[h]
\centering
\caption{Messergebnisse des invertiernden Verstärker}\label{tab:i_bode}
\begin{tabularx}{\textwidth}{YYYYY}
$f$ / \si{\hertz} & $U_A$ / \si{\volt} & $\Delta t$ / \si{\second} & $A$ / \si{\decibel} & $\varphi$ / \si{\degree} \\ \midrule
1e+01 & 4.80 & 5.00e-02 & 39.65 & 180.00 \\
1e+02 & 4.80 & 5.00e-03 & 39.65 & 180.00 \\
3e+02 & 4.80 & 1.65e-03 & 39.65 & 178.20 \\
1e+03 & 4.80 & 4.90e-04 & 39.65 & 176.40 \\
3e+03 & 4.40 & 1.52e-04 & 38.89 & 164.16 \\
1e+04 & 4.32 & 3.60e-05 & 38.73 & 129.60 \\
3e+04 & 1.30 & 9.60e-06 & 28.30 & 103.68 \\
1e+05 & 0.35 & 2.60e-06 & 16.90 & 93.60  \\
3e+05 & 0.12 & 6.80e-07 & 7.23  & 73.44  \\
8e+05 & 0.04 & 1.40e-07 & -1.94 & 40.32  \\
\end{tabularx}
\end{table}

\begin{table}[h]
\centering
\caption{Messergebnisse des nicht invertierenden Verstärker}\label{tab:ni_bode}
\begin{tabularx}{\textwidth}{YYYYY}
$f$ / \si{\hertz} & $U_A$ / \si{\volt} & $\Delta t$ / \si{\second} & $A$ / \si{\decibel} & $\varphi$ / \si{\degree} \\ \midrule
1e+01 & 5.36 & 0.00 & 40.60 & 0.00   \\
1e+02 & 5.12 & 0.00 & 40.21 & 0.00   \\
3e+02 & 5.36 & 0.00 & 40.60 & 0.00   \\
1e+03 & 5.20 & -3.20e-05 & 40.34 & -11.52  \\
3e+03 & 4.72 & -2.20e-05 & 39.50 & -23.76  \\
1e+04 & 2.80 & -1.64e-05 & 34.96 & -59.04  \\
3e+04 & 1.00 & -7.20e-06 & 26.02 & -77.76  \\
1e+05 & 0.34 & -2.50e-06 & 16.65 & -90.00  \\
3e+05 & 0.11 & -8.00e-07 & 7.16  & -86.40  \\
1e+06 & 0.03 & -3.30e-07 & -3.22 & -118.80 \\
\end{tabularx}
\end{table}

\newpage
Die in den Tabellen angeführten Daten werden in den Bodediagrammen in
\autoref{fig:bode_I} und \ref{fig:bode_NI} dargestellt.

\begin{figure}[h]
	\begin{minipage}[t]{.49\textwidth}
	\centering
	\includegraphics[height=6.4cm]{bode_I.png}
	\caption{Bodediagramm i. Verstärker}\label{fig:bode_I}
	\end{minipage}\hfill
	\begin{minipage}[t]{.49\textwidth}
	\centering
	\includegraphics[height=6.4cm]{bode_NI.png}
	\caption{Bodediagramm n.i. Verstärker}\label{fig:bode_NI}
	\end{minipage}
\end{figure}

\subsubsection{Erkenntnisse}

Anhand des Phasenverlaufs in den Bodediagrammen ist zu erkennen, dass die
Transitfrequenz wie in \cite[S.~77, Abb.~3.13]{EMTPraktikum2025} angegeben,
etwa bei $\SI{1}{\mega\hertz}$ liegt. Auch der Betrag hat etwa in diesem
Bereich seinen Null-Durchgang. Wegen des starken Rauschens und Verzerrungen
bei hohen Frequenzen kann hier jedoch nur eine ungefähre Aussage getroffen
werden. Weiters ist zu erkennen, dass die Verstärkerschaltung bei der eingestellten
Verstärkung von ca. \SI{40}{\decibel} die \SI{3}{\decibel}-Grenzfrequenz
($\Delta\varphi = \SI{45}{\degree}$) etwa bei $\SI{10}{\kilo\hertz}$ aufweist.